% Generated draft: requires mathematical and source review.
\documentclass[11pt]{article}
\usepackage[margin=25mm]{geometry}
\usepackage{fontspec}
\setmainfont{Latin Modern Roman}
\usepackage{amsmath,amssymb,mathtools,unicode-math}
\setmathfont{Latin Modern Math}
\usepackage{enumitem,needspace,xurl,hyperref}
\setlength{\parindent}{0pt}
\setlength{\parskip}{5pt}
\setlength{\emergencystretch}{3em}
\newcommand{\N}{\mathbb{N}}
\newcommand{\Z}{\mathbb{Z}}
\newcommand{\Q}{\mathbb{Q}}
\newcommand{\R}{\mathbb{R}}
\newcommand{\C}{\mathbb{C}}
\newcommand{\F}{\mathbb{F}}
\newcommand{\A}{\mathbb{A}}
\newcommand{\PP}{\mathbb{P}}
\newcommand{\E}{\mathbb{E}}
\newcommand{\eps}{\varepsilon}
\newcommand{\dd}{\,\mathrm{d}}
\begin{document}
\section*{UnsolvedMath ID: 9400004}
\textbf{Source problem code:} catalogue entry
\subsection{Definitions and mathematical statement}
The exponent pair conjecture asks whether
\[
 \forall\varepsilon>0,\qquad (\varepsilon,1/2+\varepsilon)
 \text{ is an exponent pair}.
\]
An exponent pair encodes a uniform bound for one-dimensional oscillatory sums. For clarity, one classical phase convention uses \(e(t)=\exp(2\pi i t)\) and asserts a bound of the form
\[
 \left|\sum_{a<m\leq b}e(f(m))\right|
 \leq C\bigl((ya^{-s-1})^k a^l+(ya^{-s-1})^{-1}\bigr).
\]
Here \(s>0\), \(a\geq1\), \(a\leq b\leq2a\), \(y>0\), and \(f\) belongs to a prescribed smooth class close in sufficiently many derivatives to \(x\mapsto yx^{-s}\). For example, a derivative convention requires
\[
 \left|f^{(j)}(x)-y\frac{\dd^j}{\dd x^j}x^{-s}\right|
 \leq\delta yx^{-s-j}\qquad(1\leq j\leq J,\ a\leq x\leq b),
\]
with \(J,\delta,C\) depending on fixed phase and exponent parameters but not on \(a,b,y,f\). The exact phase class and admissible-pair convention are not supplied in the excerpt, so this display gives conventional context rather than certifying the source's definition.

The proposed pairs approach \((0,1/2)\) as \(\varepsilon\downarrow0\); no uniform constant in \(\varepsilon\) and no endpoint bound are asserted. Definitions restricting exponent pairs to \(0\leq k\leq1/2\), \(1/2\leq l\leq1\) require reading the source's all-positive-\(\varepsilon\) wording through arbitrarily small positive \(\varepsilon\) or an extended convention; this needs confirmation.

The summary calls \((0,1/2)\) trivial while saying the conjectured near-endpoint pairs are unknown. Those claims are incompatible in the usual exponent-pair framework and the trivial-endpoint claim is not adopted. Its best-known and Lindel\"of-equivalence claims also require precise hypotheses and independent checking.
\subsection{Short English statement}
Can exponential-sum exponent pairs approach \((0,1/2)\) in the form \((\varepsilon,1/2+\varepsilon)\) for every arbitrarily small positive \(\varepsilon\)? The source's exponent-pair convention and the summary's inconsistent endpoint claim need review.
\subsection{Sources}
\begin{itemize}
\item Local UnsolvedMath registry, permanent ID 9400004, source code catalogue entry; no proposer or proposal year supplied; update 2026-08-06.
\item Wikipedia, \emph{List of unsolved problems in mathematics}, Number theory, General bullet 4, wikitext line 10, pinned revision 1366636767: \url{https://en.wikipedia.org/w/index.php?title=List_of_unsolved_problems_in_mathematics&oldid=1366636767#Number_theory}. Registry access date 2026-07-29; MediaWiki extraction. The provenance stages this item pending final semantic deduplication and primary-source status review, and marks redistribution terms for review. Wikipedia attribution and CC BY-SA 4.0 requirements must be retained for release.
\item Supplied summary references: Wikipedia's ``Exponent pair'' entry and Bourgain, ``On the distribution of Dirichlet sums'' (2017), with the reported pair \((13/84,55/84)\). No complete bibliography or exact epsilon qualifications are supplied. The summary's status and optimality claims are not independently verified.
\end{itemize}
\end{document}
